% Vía B, presentación del tema. Compilar desde proposal/.
\documentclass[11pt]{article}
\usepackage[utf8]{inputenc}
\usepackage[T1]{fontenc}
\usepackage[english]{babel}
\usepackage[margin=1in]{geometry}
\usepackage{amsmath,amssymb,amsthm}
\usepackage{xcolor}
\usepackage{titlesec}
\usepackage[colorlinks=true,allcolors=blue!55!black]{hyperref}
\definecolor{azul}{HTML}{123B65}
\titlespacing*{\section}{0pt}{1.3ex}{0.45ex}
\begin{document}
\begin{center}
{\large\color{azul}\bfseries Fricciones de ejecución y desempeño de la inversión pública local}\\[0.3em]
Isis Milagros Roque Tinta · Vía B: modelo económico de tesis\\[-0.1em]
{\small Presentación del tema · 14 de octubre de 2026 · \url{https://github.com/isisroquet-mi/ai-project}}
\end{center}
\vspace{-0.7em}
\section{Pregunta y vía}
La tesis de Bolo Añorga y Roque Tinta (2026), \emph{Pobreza y desempeño en la inversión pública en el Perú}, pregunta cómo se relaciona la pobreza distrital con la velocidad de ejecución y la variación de costos de proyectos de gobiernos locales. Su análisis empírico mide la velocidad mediante avance financiero acumulado por año transcurrido desde el primer devengado y la variación de costos mediante el logaritmo del cociente entre costo actualizado y monto viable. La hipótesis es asociativa y contempla características del proyecto, territorio, recursos fiscales y capacidad municipal. Este proyecto elige la \textbf{Vía B}: formular un mecanismo teórico estilizado que dé una interpretación económica a ambas dimensiones, sin sustituir la estrategia empírica ni atribuir causalidad directa a la pobreza.

\section{Línea base y supuesto que se modifica}
Se toma un proyecto de tamaño normalizado a una unidad y un gobierno local que, después de conocer las condiciones de ejecución, escoge una intensidad constante $e>0$ (avance físico equivalente por unidad de tiempo). Su duración es $T(e)=1/e$. El costo directo inevitable es $M\geq0$; mientras la obra siga abierta se paga $F>0$ por unidad de tiempo; y acelerar requiere coordinación, supervisión y recursos cuyo costo \emph{por unidad de tiempo} es $\kappa e^2/2$, con $\kappa>0$. El costo total que minimiza el municipio es
\begin{equation}\label{eq:problem}
 \min_{e>0}\ K(e;\kappa)=M+\int_0^{1/e}\left(F+\frac{\kappa e^2}{2}\right)dt
 =M+\frac{F}{e}+\frac{\kappa}{2}e.
\end{equation}
La línea base supone distritos homogéneos, $\kappa=\bar\kappa>0$. Acelerar ahorra costos de permanencia $F/e$, pero eleva el costo de intensidad $\kappa e/2$. \textbf{El supuesto que se modifica} es la homogeneidad: se permite $\kappa=\kappa(p,C)$, donde $p$ es incidencia de pobreza predeterminada y $C$ es capacidad municipal, con $\kappa_p>0$ y $\kappa_C<0$ en el dominio relevante. Son \emph{supuestos condicionales sobre fricciones asociadas al contexto}, no hechos identificados ni una relación causal estimada. $M$ y $F$ permanecen constantes en esta primera comparación para aislar el canal.

El orden de eventos es: (i) en viabilidad se registra un monto previsto $K^V$ usando una fricción de referencia $\kappa^V>0$; (ii) al ejecutar se conoce $\kappa(p,C)$ y se elige $e$; (iii) se observan duración y costo actualizado. Para comparar magnitudes sin cambios de alcance, se supone el mismo $M$ y $F$ en viabilidad y ejecución. La brecha entre $\kappa^V$ y la fricción realizada representa información o restricciones no incorporadas plenamente en el presupuesto inicial.

\section{Resultado esperado y condición de primer orden}
La condición de primer orden del problema \eqref{eq:problem} y la condición de segundo orden son
\begin{equation}\label{eq:foc}
 K_e=-\frac{F}{e^2}+\frac{\kappa(p,C)}{2}=0,
 \qquad K_{ee}=\frac{2F}{e^3}>0.
\end{equation}
Como $K(e)\to\infty$ cuando $e\downarrow0$ o $e\to\infty$, existe un único óptimo interior. En la línea base, $e_0=\sqrt{2F/\bar\kappa}$; con heterogeneidad, se espera
\begin{equation}\label{eq:sol}
 e^*(p,C)=\sqrt{\frac{2F}{\kappa(p,C)}},\qquad
 T^*(p,C)=\sqrt{\frac{\kappa(p,C)}{2F}},\qquad
 K^A(p,C)=M+\sqrt{2F\kappa(p,C)}.
\end{equation}
En particular, si $\kappa_p>0$, $\partial e^*/\partial p<0$, $\partial T^*/\partial p>0$ y $\partial K^A/\partial p>0$; si $\kappa_C<0$, se invierten los tres signos. Éstas son predicciones \emph{condicionadas al supuesto sobre $\kappa$}, no resultados de la tesis empírica. Para vincular el costo con su indicador, defínase $K^V=M+\sqrt{2F\kappa^V}>0$. Entonces
\begin{equation}\label{eq:overrun}
 R(p,C)=\log\!\left(\frac{K^A(p,C)}{K^V}\right),\qquad
 R>0\ \Longleftrightarrow\ \kappa(p,C)>\kappa^V,
 \qquad R_p=\frac{F\kappa_p}{K^A\sqrt{2F\kappa}}>0.
\end{equation}
La última derivada mantiene fijo el monto viable. Una fricción mayor \emph{eleva} $R$; un \emph{sobrecosto positivo} requiere además que la fricción realizada exceda la prevista. La velocidad $e^*=1/T^*$ es una aproximación teórica al ritmo de avance financiero por año de la tesis: sólo coinciden si el gasto sigue proporcionalmente el progreso y el proyecto se completa. Esta limitación será explícita al interpretar proyectos censurados y modificaciones del costo actualizado.

\section{Antecedentes revisados y espacio de la propuesta}
Se revisaron la literatura y el anexo metodológico de la tesis (pp.\ 3--10 y 23), el resumen y referencias del artículo publicado de Loayza, Rigolini y Calvo-González (2014), el trabajo de Espinoza y Presbitero (2022), y un artículo de optimización de tiempo y costo de proyectos (Mohammadi Atashgah \emph{et al.}, 2022). Loayza \emph{et al.} analizan capacidad de gasto \emph{municipal}, pobreza y presupuesto; su variable no es la duración y el costo de \emph{cada proyecto}. Espinoza y Presbitero estiman determinantes de retrasos en proyectos internacionales y documentan su asociación con aumentos de costo; esa evidencia no deriva la decisión de intensidad local propuesta aquí. La investigación sobre optimización tiempo--costo sí estudia formalmente ese intercambio, de modo que éste \emph{no se reivindica como nuevo}; lo específico que se busca aportar es su conexión con heterogeneidad distrital, una previsión de viabilidad y las dos medidas de la tesis peruana. Esta búsqueda inicial no prueba inexistencia de un modelo equivalente: se ampliará la revisión de versiones de trabajo, apéndices y estudios posteriores antes de sostener una contribución de originalidad fuerte.

\section{Plan de demostración, simulación y riesgo}
Se demostrará existencia y unicidad del óptimo, las comparativas estáticas de \eqref{eq:sol} y la condición exacta de sobrecosto en \eqref{eq:overrun}. Se verificará simbólicamente la FOC y las derivadas con SymPy y se simularán escenarios de $\kappa$ y $\kappa^V$ para mostrar casos con $R<0$, $R=0$ y $R>0$; las simulaciones ilustrarán la lógica, sin constituir evidencia empírica. En Lean se formalizarán primero la positividad y caracterización del óptimo y luego el orden de duración y costo ante $\kappa_2>\kappa_1>0$, sin añadir las conclusiones como hipótesis. El riesgo principal es \textbf{la correspondencia de medición}: el gasto financiero por año y el costo actualizado pueden diferir de la intensidad física y del costo teórico, sobre todo en proyectos no terminados o con cambios de alcance. Si esa traducción falla, el resultado quedará presentado como mecanismo cualitativo y se evitará equiparar sus variables a las observadas.

\smallskip\noindent\textbf{Fuentes consultadas.} Bolo Añorga y Roque Tinta (2026), \emph{Pobreza y desempeño en la inversión pública en el Perú}, versión del 27/09, pp.\ 3--10, 23. Loayza, Rigolini y Calvo-González (2014), \emph{Economics \& Politics} 26(1), 56--78, \href{https://doi.org/10.1111/ecpo.12026}{doi:10.1111/ecpo.12026}. Espinoza y Presbitero (2022), \emph{International Economics} 172, 297--310, \href{https://doi.org/10.1016/j.inteco.2021.10.002}{doi:10.1016/j.inteco.2021.10.002}. Mohammadi Atashgah \emph{et al.} (2022), \emph{Journal of Industrial and Systems Engineering} 14(3), 259--279, \href{https://www.jise.ir/article_154961_2e9576261fd03c085ad26ecae80b573.pdf}{texto}. Consulta: 13/10/2026.
\end{document}
